% SGA 7 II；PDF 3--21，止于第2节标题前。
\SourcePage{3}
\SGAfrontchapter{题名与出版资料}
《数学讲义》（Lecture Notes in Mathematics），第340卷。
主编：A. Dold（海德堡）与 B. Eckmann（苏黎世）。

\begin{center}\Large \Term{algebraic-geometry}{代数几何}中的\Term{monodromy-group}{单值群}\par\normalsize Springer-Verlag\par 柏林、海德堡、纽约，1973年\end{center}
\SourcePage{4}
AMS 学科分类（1970）：主要分类 14D05、14F20；次要分类 14C99、14G10、14M10。

ISBN 3-540-06433-8，Springer-Verlag，柏林、海德堡、纽约。

ISBN 0-387-06433-8，Springer-Verlag，纽约、海德堡、柏林。

本著作受版权保护。无论涉及全部材料还是部分材料，均保留一切权利，特别是翻译、重印、插图再利用、广播、以复印机或类似方式复制以及存储于数据库的权利。
依德国版权法第54条，为私人使用以外的用途制作复制件，须向出版者支付费用，数额由与出版者的协议确定。

版权归 Springer-Verlag（柏林、海德堡），1973年。美国国会图书馆目录卡号 73-11813。德国印刷。

胶印：Julius Beltz，Hemsbach/Bergstr。
\SourcePage{5}
\begin{center}
布尔玛丽\Term{algebraic-geometry}{代数几何}讨论班\par 1967--1969\bigskip

{\Large \Term{algebraic-geometry}{代数几何}中的\Term{monodromy-group}{单值群}}\par（SGA 7 II）\bigskip

P. Deligne 与 N. Katz 著
\end{center}
\SourcePage{6}
\SGAfrontchapter{引言}
本讨论班由 A. Grothendieck 主持的第一部分，已在同一套《数学讲义》中以第288卷出版。我们只需用到第一部分第 I 讲和第 VI 讲（尤其是第5、6节）的结果。

这第二部分的关键结果，是第 XV 讲的 \Term{picard-lefschetz-formula}{Picard--Lefschetz 公式}以及第 XVIII 讲的\Term{lefschetz-pencil}{Lefschetz 铅笔}理论。

设 $X_t$ 是依赖于参数 $t$ 的一族\Term{algebraic-variety}{代数簇}。\Term{picard-lefschetz-formula}{Picard--Lefschetz 公式}描述：当参数取值 $t_0$ 时 $X_t$ 出现一个\Term{ordinary-quadratic-singularity}{普通二次奇点}，在 $t_0$ 附近 $X_t$ 的\Term{cohomology}{上同调}如何变化。它给出 $X_t$ 与 $X_{t_0}$ 的\Term{cohomology}{上同调}之间的差异，并描述 $t$ 绕 $t_0$ 旋转时的\Term{local-scheme}{局部}\Term{monodromy}{单值性}。在\Term{transcendental}{超越}理论的框架内，该公式由 Lefschetz [2] 给出，并有简单的几何意义（第 XIV 讲 3.2）。为在任意\Term{characteristic}{特征}下建立它，我们必须搭建一套庞大的工具（第 XII、XIII、XIV 讲）；第 XIII、XIV 讲的引言概述了其主要脉络。

当取无\Term{singular}{奇异}\Term{projective-variety}{射影簇} $X$ 的\Term{hyperplane-section}{超平面截面}的一个充分\Term{generic}{一般}的\Term{pencil}{铅笔}作为族 $X_t$ 时，\Term{picard-lefschetz-formula}{Picard--Lefschetz 公式}以及其他一些几何思想，给出 $X$ 的\Term{cohomology}{上同调}与其\Term{hyperplane-section}{超平面截面} $X_t$ 的\Term{cohomology}{上同调}之间的宝贵联系。只有在 $X$ 满足“Lefschetz 母牛\Term{theorem}{定理}”（théorème de Lefschetz vache，LV）（第 XVIII 讲 5.2.2）时，这些联系才能发挥全部力量。在\Term{characteristic}{特征}0时，这一条件自动满足。\Term{dimension}{维数} $\geq3$ 时，唯一已知的证明依赖于调和积分的 \Term{hodge-theory}{Hodge 理论}（Lefschetz 对关键点的证明 [2] 13，第25页有误）；若能找到另一种可推广到\Term{characteristic}{特征} $p$ 的证明，将极有意义。
\SourcePage{7}
主要应用如下。

(a) 第 XVI 讲在\Term{equal-characteristic}{等特征} $p$ 情形下，证明一个类似于 Milnor [3] 所给公式的公式，用以计算映射 $f:\mathbb C^n\to\mathbb C$ 的孤立\Term{critical-point}{临界点}处的\Term{vanishing-cycles}{消失圈}数目。

(b) 第 XIX 讲运用 Lefschetz 的方法（[2]，第106页），特别是在任意\Term{characteristic}{特征}下证明：若 $S$ 是 $\mathbb P^3$ 中\Term{degree}{次数}不等于2、3的\Term{generic}{一般}\Term{surface}{曲面}，则 $S$ 上的每条\Term{curve}{曲线}都是 $S$ 与 $\mathbb P^3$ 中另一\Term{surface}{曲面}的\Term{complete-intersection}{完全交}。

(c) 第 XX 讲采用从 Griffiths [1] 的方法移植而来的方法，在任意\Term{characteristic}{特征}下证明：存在\Term{algebraic-cycle}{代数圈}，它们\Term{homology}{同调}\Term{equivalence}{等价}于零，但其任何倍数都不\Term{algebraic-equivalence}{代数等价}于零。

还应指出，\Term{lefschetz-pencil}{Lefschetz 铅笔}方法最近已使人能够对\Term{dimension}{维数}归纳，证明下述结果（见 [4]）：若 $\mathbb F_p$ 上的无\Term{singular}{奇异}\Term{projective-variety}{射影簇} $X$ 能连同其\Term{polarization}{极化}一起\Term{lifting}{提升}到\Term{characteristic}{特征}0，且 $p\ne2$，那么 Frobenius 在
\[
H^i(X\otimes_{\mathbb F_p}\overline{\mathbb F}_p,\mathbb Q_\ell)
\]
上\Term{action}{作用}的\Term{characteristic-polynomial}{特征多项式}属于 $\mathbb Z[t]$，且与 $\ell$ 无关（$\ell\ne p$）。

第 XXII、XXI 讲是第 XX 讲（以及\Term{characteristic}{特征}2时第 XIX 讲）的准备。在第 XXII 讲，我们研究 $\mathbb F_p$ 上\Term{algebraic-variety}{代数簇}的 $\zeta$ 函数模 $p$ 的\Term{reduced}{约化}（所用方法中 Dwork 理论起着关键\Term{action}{作用}），以及 Frobenius 在各个 $H^i$ 上\Term{action}{作用}的\Term{characteristic-polynomial}{特征多项式}的\Term{integrality}{整性}（原文作 intégrabilité）。在第 XXI 讲，我们给出\Term{generic}{一般}\Term{complete-intersection}{完全交}的\Term{cohomology}{上同调}的水平的下界。

希望先对这一理论获得几何印象的读者，我建议首先阅读第 I 讲、第 XIV 讲 3.2、第 XIX 讲第4节，以及第 XIII 讲的引言。

下面的导览图（Leitfaden）大致说明各讲之间的逻辑依赖关系。
\SGAsourcegap{引言所称下接的 Leitfaden 导览图，在原 PDF 第7--9页之间未见；现有 IAS 补源第5--7页同样从引言、文献直接接目录，未收该图。}
\SourcePage{8}
\begin{description}
\item[{[1]}] P. A. Griffiths，某些积分的周期 III（On the periods of certain integrals III），\textit{Ann. of Math.}
\item[{[2]}] S. Lefschetz，位置分析与\Term{algebraic-geometry}{代数几何}（L'analysis situs et la géométrie algébrique），巴黎，1924，Gauthier-Villars；重印于 \textit{Selected Papers}（Chelsea）。
\item[{[3]}] J. Milnor，复\Term{hypersurface}{超曲面}的\Term{singular-point}{奇点}（Singular points of complex hypersurfaces），\textit{Annals of Math. Studies} 61，Princeton，1968。
\item[{[4]}] J. L. Verdier（据 P. Deligne），$\ell$ 进\Term{cohomology}{上同调}（Cohomologie $\ell$-adique），Bourbaki 讨论班，1972年11月。
\end{description}
\begin{flushright}Bures-sur-Yvette，1972年11月8日\par P. Deligne\end{flushright}
\SourcePage{9}
\SGAfrontchapter{原书目录}
以下页码为原书印刷页码；译本目录由分册统一生成。
\begin{description}
\item[第 X 讲] \Term{regular}{正则}\Term{surface}{曲面}上的交（P. Deligne），1。
\item[第 XI 讲] \Term{complete-intersection}{完全交}的\Term{cohomology}{上同调}（P. Deligne），39。
\item[第 XII 讲] \Term{quadric}{二次曲面}（P. Deligne），62。
\item[第 XIII 讲] \Term{vanishing-cycles}{消失圈}的\Term{form}{形式}体系（P. Deligne），82。
\item[第 XIV 讲] 与\Term{transcendental}{超越}理论的比较（P. Deligne），116。
\item[第 XV 讲] \Term{picard-lefschetz-formula}{Picard--Lefschetz 公式}（P. Deligne），165。
\item[第 XVI 讲] Milnor 公式（P. Deligne），197。
\item[第 XVII 讲] \Term{lefschetz-pencil}{Lefschetz 铅笔}：\Term{existence-theorem}{存在定理}（N. Katz），212。
\item[第 XVIII 讲] \Term{lefschetz-pencil}{Lefschetz 铅笔}的\Term{cohomology}{上同调}研究（N. Katz），254。
\item[第 XIX 讲] Noether \Term{theorem}{定理}（P. Deligne），328。
\item[第 XX 讲] Griffiths \Term{theorem}{定理}（N. Katz），341。
\end{description}
\SourcePage{10}
\begin{description}
\item[第 XXI 讲] \Term{complete-intersection}{完全交}的\Term{cohomology}{上同调}的水平（N. Katz），363。
\item[第 XXII 讲] 函数\SGAsourcegap{第 XXII 讲原目录题名在“fonction”之后空缺；现有补源目录亦同，未据正文题名补写。}的一个同余公式（N. Katz），401。
\end{description}
\SourcePage{11}
\SGAchapter{X}{正则曲面上的交}{Intersections sur les surfaces régulières}
\begin{center}P. Deligne\end{center}
本讲第一部分先介绍一些\Term{intersection-theory}{交理论}的预备知识（包括在 [5] 与 [7] 中），然后给出 Artin 对 Raynaud [9] 的一个\Term{theorem}{定理}的特殊情形的证明（1.13），并给出 Néron [8] 曾逐个情形核实的一个\Term{theorem}{定理}的先验性证明（1.15）。由（1.13）可相当容易地\Term{pushout}{推出} [9] 中在 [2] 中使用、从而也在第 IX 讲中使用的那些\Term{theorem}{定理}。

第二部分，我顺势继续，讲述“意大利学派”[1]、[3] 如何计算\Term{surface}{曲面}上的\Term{intersection-number}{交数}。主要结果2.9和2.11也见于 Northcott [12]、[13]。最后在2.27、2.30、2.34中，给出一个猜想关系，联系一条\Term{reduced}{既约}\Term{curve}{曲线}的\Term{singular-point}{奇点}的\Term{versal-formal-moduli-scheme}{通用形式模概形}（schéma formel versel des modules）的\Term{dimension}{维数}、\Term{singular-point}{奇点}的\Term{invariants}{不变量} $\delta$，以及一个度量其\Term{symmetry}{对称性}的\Term{invariants}{不变量}。

本讲不会在讨论班后续内容中使用。
\SGAsection{1}{算术曲面上的交数}
\textbf{1.0。} 本节给定一个\Term{dimension}{维数}2的\Term{regular-scheme}{正则概形} $X$、$X$ 的一个\Term{reduced}{既约}\Term{divisor}{除子} $Y$，以及 $H^0(Y,\mathcal O_Y)$ 的一个子域 $k$。假定 $Y$ 是 $k$ 上的\Term{proper}{固有}\Term{curve}{曲线}（EGA II 7.4.1）。下面给出这种情形的三个例子。
\begin{SGAresult}{例}{1.1}
设 $(S,\eta,s)$ 是一个\Term{trait}{离散赋值环谱}，$X$ 是一个 $S$-\Term{scheme}{概形} $f:X\to S$，它\Term{regular}{正则}，且在 $S$ 上\Term{proper}{固有}、\Term{flat}{平坦}，所有\Term{fibre}{纤维}纯为一维。取 $Y=(f^{-1}(s))_{\mathrm{red}}$，$k=k(s)$。从1.10起，我们将只考虑这一情形。
\end{SGAresult}
\SourcePage{12}
\begin{SGAresult}{例}{1.2}
设 $X_0$ 是一个\Term{dimension}{维数}2的\Term{normal}{正规}\Term{local-ring}{局部环}的谱，$s$ 是 $X_0$ 的\Term{closed-point}{闭点}，$f:X\to X_0$ 是一个\Term{regular}{正则}、\Term{connected}{连通}的 $X_0$-\Term{scheme}{概形}，它在 $X_0$ 上\Term{proper}{固有}，不同于 $X_0$，但在 $X_0\setminus\{s\}$ 上与 $X_0$ 一致。取 $Y=(f^{-1}(s))_{\mathrm{red}}$，$k=k(s)$。
\end{SGAresult}
\begin{SGAresult}{例}{1.3}
取 $X$ 为域 $k$ 上的\Term{complete}{完备}\Term{regular}{正则}\Term{surface}{曲面}，$Y$ 为 $X$ 上的一条\Term{reduced}{既约}\Term{curve}{曲线}。
\end{SGAresult}
回顾下述\Term{lemma}{引理}（SGA 6 IV，练习1.9）。
\begin{SGAresult}{\Term{lemma}{引理}}{1.4}
设 $Z_0$ 是一个\Term{noetherian}{诺特}\Term{scheme}{概形} $Z$ 的\Term{closed-subscheme}{闭子概形}。则 $Z_0$ 到 $Z$ 的包含，诱导一个\Term{isomorphism}{同构}：从\Term{coherent-sheaf}{凝聚} $\mathcal O_{Z_0}$-模\Term{sheaf}{层}\Term{category}{范畴}的 Grothendieck 群，到 $Z$ 上\Term{support}{支集}在 $Z_0$ 中的\Term{coherent-sheaf}{凝聚层}\Term{category}{范畴}的 Grothendieck 群。
\end{SGAresult}
这个\Term{lemma}{引理}使我们能为 $X$ 上\Term{support}{支集}在 $Y$ 中的\Term{coherent-sheaf}{凝聚层} $\mathcal F$ 定义 \Term{euler-poincare-characteristic}{Euler--Poincaré 特征} $\chi(\mathcal F)$，使它满足：

(i) 若 $\mathcal F$ 是一个\Term{coherent-sheaf}{凝聚} $\mathcal O_Y$-模\Term{sheaf}{层}，则
\begin{equation}\chi(\mathcal F)=\dim_k H^0(\mathcal F)-\dim_k H^1(\mathcal F).\tag{1.4.1}\end{equation}
(ii) $\chi$ 对\Term{short-exact-sequence}{短正合列}可加。

若 $\mathcal O_X$-模\Term{complex}{复形} $K$ 的\Term{cohomology-sheaf}{上同调层}是\Term{support}{支集}在 $Y$ 中的\Term{coherent-sheaf}{凝聚层}，置
\begin{equation}\chi(K)=\sum_i(-1)^i\chi(\mathcal H^i(K)).\tag{1.4.2}\end{equation}
若 $\mathcal F$ 集中于 $Y$ 的一个点 $y$，则
\begin{equation}\chi(\mathcal F)=\operatorname{lg}_{\mathcal O_y}(\mathcal F_y)\,[k(y):k].\tag{1.4.3}\end{equation}
\SourcePage{13}
在（1.1）、（1.2）或（1.3）的假设下，有
\begin{equation}\chi(\mathcal F)=\operatorname{lg}f_*\mathcal F-\operatorname{lg}R^1f_*\mathcal F.\tag{1.4.4}\end{equation}
设 $D$ 是一个集中在 $Y$ 上的\Term{effective-divisor}{有效除子}，$\mathcal L$ 是 $D$ 上的一个\Term{invertible-sheaf}{可逆层}。则有且只有一种方式定义\Term{degree}{次数} $\deg_D(\mathcal L)$，使得
\begin{equation}
\begin{cases}
\deg_D(\mathcal L'\otimes\mathcal L'')=\deg_D(\mathcal L')+\deg_D(\mathcal L''),\\
\text{若 }\mathcal L\subset\mathcal O,\quad\deg_D(\mathcal L)=-\chi(\mathcal O/\mathcal L).
\end{cases}\tag{1.4.5}
\end{equation}
于是有 $D$ 上的 \Term{riemann-roch-formula}{Riemann--Roch 公式}
\begin{equation}\chi(\mathcal L)=\chi(\mathcal O_D)+\deg_D(\mathcal L).\tag{1.4.6}\end{equation}
设 $D=\sum_i n_iD_i$ 是 $D$ \Term{resolution}{分解}为\Term{reduced}{既约}\Term{irreducible}{不可约}\Term{divisor}{除子}的\Term{resolution}{分解}；有
\begin{equation}\deg_D(\mathcal L)=\sum_i n_i\deg_{D_i}(\mathcal L).\tag{1.4.7}\end{equation}
\begin{SGAresult}{定义}{1.5}
设 $D,E$ 是 $X$ 上两个\Term{effective-divisor}{有效除子}，且 $D$ 集中在 $Y$ 上。定义 $D,E$ 的\Term{intersection-number}{交数}（相对于域 $k$（1.0））为
\[
(D,E)=\chi(\mathcal O_D\mathbin{\overset{\mathrm L}{\otimes}}\mathcal O_E)
=\chi(\mathcal O_D\otimes\mathcal O_E)-\chi(\operatorname{Tor}^1(\mathcal O_D,\mathcal O_E)).
\]
\end{SGAresult}
事实上，在所作假设下有 $\operatorname{Tor}^2(\mathcal O_D,\mathcal O_E)=0$。
\begin{SGAresult}{命题}{1.6}
在1.5的假设下，有：

(i) \Term{intersection-number}{交数} $(D,E)$ 对\Term{divisor}{除子} $D$ 与 $E$ 都可加；由此可通过\Term{linear}{线性}延拓，定义两个不一定\Term{effective-divisor}{有效}的\Term{divisor}{除子}的\Term{intersection-number}{交数}。
\SourcePage{14}

(ii) $(D,E)=\deg_D(\mathcal O(E))$。

(iii) 若 $D,E$ 只在\Term{finite}{有限}多个点相交，则 $(D,E)=\chi(\mathcal O_D\otimes\mathcal O_E)\geq0$；等号成立仅当 $D,E$ \Term{disjoint}{不相交}。

(iv) 若 $D,E$ 都集中在 $Y$ 上，则 $(D,E)=(E,D)$。

(v) 若 $N$ 是 $D$ 的\Term{conormal-sheaf}{余法层}，则 $(D)^2=(D,D)=-\deg_D(N)$。

(vi) 若 $E$ 是 $X$ 上一个\Term{meromorphic}{亚纯}函数的\Term{divisor}{除子}（不一定\Term{effective-divisor}{有效}），则 $(D,E)=0$。
\end{SGAresult}
\begin{proof}
\Term{short-exact-sequence}{短正合列}
\[
0\longrightarrow\mathcal O(-E)\longrightarrow\mathcal O_X\longrightarrow\mathcal O_E\longrightarrow0
\]
结合（1.4.6）给出
\begin{align*}
(D,E)&\overset{\mathrm{def}}=\chi(\mathcal O_D\mathbin{\overset{\mathrm L}{\otimes}}\mathcal O_E)
=\chi(\mathcal O_D\mathbin{\overset{\mathrm L}{\otimes}}\mathcal O_X)-\chi(\mathcal O_D\mathbin{\overset{\mathrm L}{\otimes}}\mathcal O(-E))\\
&=\chi(\mathcal O_D)-\chi(\mathcal O(-E)|_D)
=-\deg_D(\mathcal O(-E))=\deg_D(\mathcal O(E)).
\end{align*}
这证明了(ii)、(vi)，并由（1.4.5）、（1.4.7）证明了双可\Term{additive}{加性}。在(iii)的假设下，$\operatorname{Tor}^1(\mathcal O_D,\mathcal O_E)=0$，从而得到结论。在(v)的假设下，有
\[
\mathcal O_D\otimes\mathcal O_D=\mathcal O_D,
\qquad\operatorname{Tor}^1(\mathcal O_D,\mathcal O_D)\simeq N
\quad\text{（SGA 6 VII 2.5）},
\]
故
\[
(D,D)=\chi(\mathcal O_D)-\chi(N)=-\deg_D(N).
\]
最后，(iv)由\Term{tensor-product}{张量积}的\Term{symmetry}{对称性}得到。
\end{proof}
应注意，(ii)、(iv)、(vi)在不假定\Term{divisor}{除子}\Term{effective-divisor}{有效}时仍成立。
\SourcePage{15}
\begin{SGAresult}{命题}{1.7}
设 $(Y_i)_{1\leq i\leq k}$ 是 $Y$ 的\Term{irreducible-component}{不可约分支}，$f$ 是 $X$ 上的\Term{meromorphic}{亚纯}函数，且
\[
\operatorname{div}(f)=\sum_i n_iY_i+Z
\]
为其\Term{divisor}{除子}，其中 $|Z|\cap Y$ 为\Term{finite}{有限}集。则置 $D_i=n_iY_i$，对于 $(a_i)\in\mathbb Q^k$ 有
\[
\Bigl(\sum_i a_iD_i\Bigr)^2
=-\sum_{i<j}(a_i-a_j)^2(D_i,D_j)-\sum_i a_i^2(D_i,Z).
\]
\end{SGAresult}
\begin{proof}[证明（据 [7]）]
由1.6 (v)、(vi)，有
\begin{align*}
\Bigl(\sum_i a_iD_i\Bigr)^2
&=\sum_i a_i\Bigl(D_i,\sum_j a_jD_j\Bigr)
=\sum_i a_i\Bigl(D_i,\sum_j a_jD_j-a_i\operatorname{div}f\Bigr)\\
&=\sum_i a_i\Bigl(D_i,\sum_{j\ne i}(a_j-a_i)D_j-a_iZ\Bigr)\\
&=\sum_{i\ne j}a_i(a_j-a_i)(D_i,D_j)-\sum_i a_i^2(D_i,Z)\\
&=-\sum_{i<j}(a_i-a_j)^2(D_i,D_j)-\sum_i a_i^2(D_i,Z).
\end{align*}
\end{proof}
\begin{SGAresult}{推论}{1.8}
在例1.1的假设下，取 $f$ 为 $S$ 的一个\Term{uniformizer}{一致化参数} $\pi$。此时 $Z=0$，且
\[
\Bigl(\sum_i k_iY_i\Bigr)^2
=-\sum_{i<j}\frac{(k_in_j-k_jn_i)^2}{n_in_j}(Y_i,Y_j).
\]
因此，\Term{intersection-form}{交形式}为半负定；若 $Y$ \Term{connected}{连通}，则其核只包含 $\operatorname{div}(\pi)$ 的倍数。
\end{SGAresult}
\begin{SGAresult}{推论}{1.9（Mumford [7]）}
在例1.2的假设下，取 $f$ 为\Term{local-ring}{局部环} $\mathcal O_s$ 的\Term{maximal-ideal}{极大理想}中的非零元素。此时 $Z>0$，因此\Term{intersection-form}{交形式}为负定。
\end{SGAresult}
\SourcePage{16}
\textbf{1.10。} 下文只考虑例（1.1）中的情形。

Lichtenbaum [5] 已证明，在这种情形下 $X$ 自动是 $S$ 上的\Term{projective}{射影}的，因而定义\Term{complex}{复形} $Rf^!\mathcal O_S$（[4] III 8.7，第190页）并不费力。由于 $f$ 的\Term{relative-dimension}{相对维数}为1，且 $X,S$ 均\Term{regular}{正则}（Gorenstein 已足够），存在 $S$ 上的一个\Term{invertible-sheaf}{可逆层} $\omega_{X/S}$，在唯一\Term{isomorphism}{同构}的意义下唯一（[4] V 8.3、V 9.1），使得
\begin{equation}Rf^!\mathcal O_S\simeq\omega_{X/S}[1].\tag{1.10.1}\end{equation}
用 $K$ 表示任意一个满足 $\omega\simeq\mathcal O(K)$ 的\Term{divisor}{除子}（\Term{effective-divisor}{有效}与否均可）。

设 $D$ 是一个集中在 $Y$ 上的\Term{effective-divisor}{有效除子}，$N$ 是其\Term{conormal-sheaf}{余法层}：
\[
\begin{tikzcd}
D\arrow[r,"i"]\arrow[dr,"g"']&X\arrow[d,"f"]\\
&S.
\end{tikzcd}
\]
有
\begin{equation}
Rg^!\mathcal O_S=Ri^!Rf^!\mathcal O_S=Ri^!\omega_{X/S}[1]
=\omega_{X/S}|_D\otimes N^\vee[0].\tag{1.10.2}
\end{equation}
由\Term{duality-theorem}{对偶定理}（[4] V 11.1，第210页），
\begin{equation}
Rg_*Rg^!\mathcal O_S\simeq R\operatorname{Hom}(Rg_*\mathcal O_D,\mathcal O_S).\tag{1.10.3}
\end{equation}
此外，对于任何集中在\Term{closed-point}{闭点}的 $\mathcal O_S$-模 $M$，有
\begin{equation}\chi(R\operatorname{Hom}(M,\mathcal O_S))=-\chi(M).\tag{1.10.4}\end{equation}
因此，在此利用（1.4.4）和（1.10.3），得到
\[
\chi(Rg^!\mathcal O_S)=-\chi(\mathcal O_D).
\]
又由（1.4.6），
\[
\chi(Rg^!\mathcal O_S)=\chi(\mathcal O_D)+\deg_D(Rg^!\mathcal O_S),
\]
\SourcePage{17}
结合（1.10.2）得
\[
-2\chi(\mathcal O_D)=\deg_D(Rg^!\mathcal O_S)
=\deg_D(\omega|_D\otimes N^\vee)=\deg_D(\omega)-\deg_D(N),
\]
依1.6 (ii)、(v)，这可改写为：
\begin{SGAresult}{命题}{1.11}
在（1.1）的假设下，对于任何集中在 $Y$ 上的\Term{effective-divisor}{有效除子} $D$，有
\[
\chi(\mathcal O_D)=-\frac12(D,D+K).
\]
\end{SGAresult}
\textbf{1.12。} 用 $d$ 表示 $X_s=f^{-1}(s)$ 的各\Term{irreducible-component}{不可约分支}的几何\Term{multiplicity}{重数}的最大公因数，即 $f$ 的几何特殊\Term{fibre}{纤维}的\Term{maximal-point}{极大点}的\Term{local-ring}{局部环}\Term{length}{长度}的最大公因数。若 $X_\eta$ 有一个\Term{rational-point}{有理点}，即 $f:X\to S$ 有一个\Term{section}{截面} $x$，则 $d=1$；更确切地说，$f$ 在 $x(s)$ 处\Term{smooth}{光滑}。事实上，$x$ 是\Term{regular-scheme}{正则概形}之间的\Term{closed-immersion}{闭浸入}，所以是\Term{regular-immersion}{正则浸入}，然后应用 EGA IV 17.12.1。

下面给出 M. Artin 对 M. Raynaud [9] 的一个\Term{theorem}{定理}的特殊情形的证明；Raynaud 的\Term{theorem}{定理}只假定 $d$ 与 $k$ 的\Term{characteristic}{特征} $p$ 互素。
\begin{SGAresult}{\Term{theorem}{定理}}{1.13}
沿用上述记号，若 $Y$ \Term{geometrically-connected}{几何连通}，且 $d=1$，则 $X$ 在 $S$ 上于\Term{dimension}{维数}0\Term{cohomologically-flat}{上同调平坦}（EGA III 7.8.1）；更精确地说（EGA III 7.6.9 (ii)），有
\[
k\xrightarrow{\sim}H^0(X_s,\mathcal O_{X_s}).
\]
\end{SGAresult}
\begin{proof}
容易归约到 $S$ \Term{strictly-local}{严格局部}的情形。由于 $Y$ \Term{geometrically-connected}{几何连通}，$H^0(Y,\mathcal O_Y)$ 是 $k$ 的\Term{radicial-extension}{根性扩张}；又由于 $(d,p)=1$，$Y$ 至少有一个\Term{irreducible-component}{不可约分支}是\Term{geometrically-reduced}{几何既约}的，\SourcePage{18}由此容易\Term{pushout}{推出} $H^0(Y,\mathcal O_Y)=k$。

设 $D$ 是 $X$ 上的一个\Term{divisor}{除子}，满足
\[
Y\leq D<X_s.
\]
由于 $d=1$，$D$ 不是 $X_s$ 的有理数倍，故由1.8，
\[
(D,X_s-D)=-(D,D)>0.
\]
所以\Term{divisor}{除子} $Y$ 有一个\Term{irreducible-component}{不可约分支} $E$，使得
\[
0<E\leq X_s-D,\qquad(D,E)>0.
\]
考虑\Term{commutative-diagram}{交换图}
\[
\begin{tikzcd}[column sep=small]
0\arrow[r]&\mathcal O(-D-E)\arrow[r]\arrow[d,"\varphi"]&\mathcal O\arrow[r]\arrow[d,equal]&\mathcal O_{D+E}\arrow[r]\arrow[d,"\psi"]&0\\
0\arrow[r]&\mathcal O(-D)\arrow[r]&\mathcal O\arrow[r]&\mathcal O_D\arrow[r]&0.
\end{tikzcd}
\]
\Term{snake-lemma}{蛇引理}给出 $\operatorname{Ker}(\psi)$ 与 $\operatorname{Coker}(\varphi)$ 之间的\Term{isomorphism}{同构}。这个\Term{cokernel}{余核}正是 $\mathcal O(-D)$ 在 $E$ 上的\Term{restriction}{限制}，因而有\Term{short-exact-sequence}{短正合列}
\[
0\longrightarrow\mathcal O(-D)|_E\longrightarrow\mathcal O_{D+E}\longrightarrow\mathcal O_D\longrightarrow0.
\]
依构造，$\deg_E(\mathcal O(-D))=-(D,E)<0$，故 $H^0(\mathcal O(-D)|_E)=0$，而\Term{restriction}{限制}\Term{arrow}{箭头}
\[
0\longrightarrow H^0(\mathcal O_{D+E})\longrightarrow H^0(\mathcal O_D)
\]
是\Term{injectivity}{单射}。由此可归纳构造一列\Term{divisor}{除子}
\[
Y=D_0\leq D_1\leq\cdots\leq D_n=X_s,
\]
使 $H^0(\mathcal O_{D_i})$ \Term{injectivity}{单射}到 $H^0(\mathcal O_{D_{i-1}})$ 中。因此 $H^0(\mathcal O_{X_s})=H^0(\mathcal O_Y)=k$，这就完成了1.13的证明。
\end{proof}
\SourcePage{19}
下述事实的证明见 [5] 或 [10]。
\begin{SGAresult}{\Term{theorem}{定理}}{1.14（Castelnuovo、Enriques 等）}
在（1.1）的假设下，设 $D$ 是包含于 $Y$ 的\Term{reduced}{既约}\Term{irreducible}{不可约}\Term{divisor}{除子}。假定 $H^1(D,\mathcal O_D)=0$，且置 $k'=H^0(D,\mathcal O_D)$ 后，有
\[
(D,D)=-[k':k].
\]
则 $D$ 是 $k'$ 上的\Term{projective-line}{射影直线}，且存在一个\Term{regular}{正则}的 $S$-\Term{scheme}{概形} $X_1$、$X_1$ 的一个\Term{residue-field}{剩余域}为 $k'$ 的\Term{closed-point}{闭点} $x_1$，以及 $X$ 与从 $X_1$ 沿 $x_1$ \Term{blow-up}{爆破}所得的 $S$-\Term{scheme}{概形}之间的\Term{isomorphism}{同构}，这个\Term{isomorphism}{同构}将 $D$ 与 $x_1$ 在\Term{blow-up}{爆破}中的\Term{fibre}{纤维}相等同。
\end{SGAresult}
具有这些性质的\Term{divisor}{除子} $D$ 称为\Term{exceptional-curve-first-kind}{第一类例外曲线}。

我忍不住要证明 Néron [8] 命题 $1'$（第90页）的下述命题。
\begin{SGAresult}{命题}{1.15}
假定 $X_\eta$ 是一条\Term{elliptic-curve}{椭圆曲线}（特别地，它有一个\Term{section}{截面}），且 $X$ 不包含任何\Term{exceptional-curve-first-kind}{第一类例外曲线}。则用（1.10.1）的记号，有
\[
f^*f_*\omega_{X/S}\xrightarrow{\sim}\omega_{X/S},
\]
且\Term{universal}{普遍}成立。
\end{SGAresult}
\begin{proof}
由于 $d=1$（1.12），依（1.13），只需证明 $\omega_{X/S}$ \Term{isomorphism}{同构}于 $\mathcal O_X$。由于 $\omega_{X/S}$ 在\Term{elliptic-curve}{椭圆曲线} $X_\eta$ 上\Term{trivial}{平凡}，事先就知道
\[
\omega_{X/S}\simeq\mathcal O\Bigl(\sum_i n_iY_i\Bigr),
\]
\SourcePage{20}
其中 $Y_i$ 为 $Y$ 的\Term{irreducible-component}{不可约分支}；又由于 $\operatorname{div}(\pi)\sim0$，可假定各 $n_i$ 非负。若 $\sum_i n_iY_i$ 与\Term{divisor}{除子} $\operatorname{div}(\pi)$ 成比例，则它是 $\operatorname{div}(\pi)$ 的整数倍，结论已得。否则，有
\[
\Bigl(\sum_i n_iY_i\Bigr)^2<0,
\]
存在 $i$，使
\[
\Bigl(Y_i,\sum_jn_jY_j\Bigr)<0,
\]
然后应用下面的\Term{lemma}{引理}。
\end{proof}
\begin{SGAresult}{\Term{lemma}{引理}}{1.16}
在1.1的假设下，假定 $Y$ \Term{connected}{连通}，且 $Y$ \Term{reducible}{可约}，或者 $\chi(\mathcal O_Y)\leq0$。则 $Y$ 的一个\Term{irreducible-component}{不可约分支} $Y_i$ 是\Term{exceptional-curve-first-kind}{第一类例外曲线}的充要条件是 $(Y_i,K)<0$。
\end{SGAresult}
\begin{proof}
置 $k'=H^0(Y_i,\mathcal O_{Y_i})$。\Term{riemann-roch-formula}{Riemann--Roch 公式}1.11给出
\[
\chi(\mathcal O_{Y_i})=[k':k]\bigl(1-\dim_{k'}H^1(Y_i,\mathcal O_{Y_i})\bigr)
=-\frac12\bigl((Y_i,Y_i)+(Y_i,K)\bigr).
\]
若 $Y_i$ 是\Term{exceptional-curve-first-kind}{第一类例外曲线}，因而有 $(Y_i,K)=-[k':k]<0$。

反之，若 $(Y_i,K)<0$，则 $\chi(\mathcal O_{Y_i})>0$，这意味着 $\chi(\mathcal O_{Y_i})=[k':k]$，且 $Y$ \Term{reducible}{可约}，从而由1.8得 $(Y_i,Y_i)<0$。负数 $(Y_i,Y_i)$ 与 $(Y_i,K)$ 都被 $[k':k]$ 整除，因此必有
\[
(Y_i,Y_i)=(Y_i,K)=-[k':k].
\]
\end{proof}
\textbf{练习 (a)。} 在（1.1）的假设下，设 $X_\eta$ 是一条\Term{elliptic-curve}{椭圆曲线}下的\Term{principal-homogeneous-space}{主齐性空间}，且 $X$ 不包含任何\Term{exceptional-curve-first-kind}{第一类例外曲线}。证明在 $X$ 上，\Term{riemann-roch-formula}{Riemann--Roch 公式}（1.11）写作
\[
\chi(\mathcal O_D)=-\frac12(D,D).
\]
\SourcePage{21}
\textbf{(b)} 再假定\Term{residue-field}{剩余域} $k(s)$ \Term{algebraically-closed}{代数闭}。证明 $X_s$ 的\Term{irreducible-component}{不可约分支}图是 Néron [8] 第124--125页第四栏所考虑的某个图（情形 a 至 c8）的倍数。这个图确定各\Term{reduced}{既约}\Term{irreducible-component}{不可约分支}的种数、它们在 $X_s$ 中的\Term{multiplicity}{重数}，以及 $(X_s)_{\mathrm{red}}$ 的\Term{singular-point}{奇点}的类型和配置。

先处理 $X_s$ \Term{irreducible}{不可约}的情形；排除这一情形后，证明每一个\Term{reduced}{既约}\Term{irreducible-component}{不可约分支}都是\Term{self-intersection}{自交数}为 $-2$ 的\Term{projective-line}{射影直线}。若 $X_s=\sum_i n_iD_i$ 是 $X_s$ \Term{resolution}{分解}为\Term{reduced}{既约}\Term{irreducible}{不可约}\Term{divisor}{除子}的\Term{resolution}{分解}，则将这个条件改写为
\begin{equation}n_i=\frac12\sum_{j\ne i}n_i(D_iD_j).\tag{*}\end{equation}
另一方面，$X_s$ \Term{connected}{连通}，所以：
\begin{equation}
\begin{gathered}
\text{对于任意 }i,j\text{，存在一条链 }i=i_1,i_2,\ldots,i_n=j,\\
\text{使得 }(D_{i_k}D_{i_{k+1}})\ne0.
\end{gathered}\tag{**}
\end{equation}
然后求（*）与（**）的所有整数解。先证明：除 Néron 的情形 b2、c2 外，\Term{intersection-number}{交数} $(D_iD_j)$ 等于0或1。
